\worksheetheader{COMP 469}{Fall 2026}{\ifsolutions Discussion 2 Solutions\else Regular Discussion 2\fi}

\begin{center}
{\large\textbf{Intelligent Agents} \quad (AIMA 4th ed., Chapter 2)}
\end{center}

\vspace{0.5em}

\instrnote{Budget 50--60 minutes. Suggested pacing: Q1 (8 min), Q2 (15 min),
Q3 (15 min), Q4 (8 min), Q5 (12 min). Q2(c) and Q3(b) are the two parts worth reserving
whiteboard time for; both reward argument over recall. Common misconceptions to watch for are
flagged in the individual solutions below.}

% =====================================================================
\qhead{Q1. The Vacuum World, Formalized}

Consider the vacuum world of Figure 2.2, generalized to a $1 \times n$ row of squares
($n \geq 2$). Each square is either \textit{Clean} or \textit{Dirty}. The agent occupies
exactly one square at a time. Its sensors report the label of its current square and whether
that square contains dirt. Its actions are $\{\textit{Left}, \textit{Right}, \textit{Suck},
\textit{NoOp}\}$. Sucking cleans the current square, clean squares stay clean, and a move
that would leave the row leaves the agent where it is.

\begin{center}
\begin{tikzpicture}[scale=1.0]
  \draw[thick] (0,0) rectangle (2,2);
  \draw[thick] (2,0) rectangle (4,2);
  \node at (0.3,1.7) {\textbf{A}};
  \node at (2.3,1.7) {\textbf{B}};
  % dirt in B
  \foreach \p in {(2.6,0.5),(3.0,0.7),(3.4,0.4),(3.1,0.3),(2.9,0.55)}
    {\fill[gray] \p circle (0.06);}
  % agent in A
  \draw[thick,fill=gray!25] (0.7,0.4) rectangle (1.5,1.1);
  \draw[thick] (1.1,1.1) -- (1.1,1.45);
  \draw[thick,fill=gray!40] (1.1,1.55) circle (0.12);
\end{tikzpicture}
\end{center}

\begin{enumerate}
\item State precisely the difference between an \textbf{agent function} and an
\textbf{agent program}. Which of the two can depend on the entire percept history, and what
must the other one do to obtain the same behavior?
\ansspace{4.0cm}
\begin{sol}
The \textbf{agent function} is an abstract mathematical mapping from every possible
\emph{percept sequence} to an action. It is an external characterization of behavior; it does
not say how the behavior is produced.

The \textbf{agent program} is a concrete implementation running on an architecture
($\textit{agent} = \textit{architecture} + \textit{program}$). It takes only the
\emph{current percept} as input, because that is all the sensors deliver at that instant.

Only the agent function may depend on the whole history. If the desired behavior depends on
the history, the program has to store what it needs in \textbf{internal state} across calls.
That requirement is exactly what separates a simple reflex agent from a model-based one, and
it is the crux of Q2(c).
\end{sol}

\item How many distinct \textbf{environment states} are there? Justify your count.
\ansspace{2.5cm}
\begin{sol}
$n \cdot 2^{n}$. The state is (agent location, dirt configuration): $n$ choices of square for
the agent, and each of the $n$ squares independently Clean or Dirty, giving $2^{n}$ dirt
configurations. For the two-square world of Figure 2.2 this is $2 \cdot 2^{2} = 8$ states.

Common error: answering $2^{n}$ (forgetting the agent's own position is part of the state) or
$n + 2^n$ (adding instead of multiplying independent factors).
\end{sol}

\item How many distinct percepts can the agent receive? How many distinct percept sequences of
length at most $T$ are there? Give a closed form.
\ansspace{3.0cm}
\begin{sol}
A percept is (location, dirt status of that square), so there are $|P| = 2n$ distinct percepts.

The number of percept sequences of length exactly $t$ is $(2n)^{t}$, so the number of sequences
of length at most $T$ is
\[
\sum_{t=1}^{T} (2n)^{t} \;=\; \frac{(2n)^{T+1} - 2n}{2n - 1}.
\]
This is precisely the number of rows a \textsc{Table-Driven-Agent} would need
(AIMA \S 2.4.1). For $n = 2$ and a modest lifetime of $T = 20$ that is already about
$1.5 \times 10^{12}$ rows.
\end{sol}

\item How many distinct agent functions are there over percept sequences of length at most $T$?
Evaluate your expression for $n = 2$, $T = 2$, and say in one sentence what the number implies
about building agents by table lookup.
\ansspace{3.0cm}
\begin{sol}
Each of the $N = \sum_{t=1}^{T}(2n)^{t}$ sequences may be mapped independently to any of the
$|A| = 4$ actions, so there are $|A|^{N} = 4^{N}$ agent functions.

For $n = 2$, $T = 2$: $N = 4^{1} + 4^{2} = 20$, so $4^{20} \approx 1.1 \times 10^{12}$ agent
functions --- for a two-square world with a two-step lifetime.

Implication: the space of behaviors is astronomically larger than the space of \emph{sensible}
behaviors, and tabulating it is hopeless. Even a correct table is (a) unstorable, (b)
unwritable by a designer, and (c) unlearnable from experience. The design problem is to
produce rational behavior from a small program, not from a vast table.
\end{sol}
\end{enumerate}

\newpage
% =====================================================================
\qhead{Q2. Rationality Is Relative}

Work in the two-square world (squares A and B). Consider the agent whose agent function is:
\textit{if the current square is dirty then Suck; otherwise move to the other square.}
Assume the geography is known a priori, but the initial dirt distribution and initial location
are not; sensors are accurate; clean squares stay clean.

\begin{enumerate}
\item Under performance measure $M_1$ --- one point for each clean square at each time step over
a lifetime of 1000 steps --- is this agent rational? Justify against the four things rationality
depends on.
\ansspace{4.0cm}
\begin{sol}
Yes. Rationality depends on (1) the performance measure, (2) the agent's prior knowledge of the
environment, (3) the available actions, and (4) the percept sequence to date. Under $M_1$:

\begin{itemize}[leftmargin=1.5em]
\item Dirt is only removable by \textit{Suck}, and the agent sucks whenever it perceives dirt,
so it never delays a cleaning it could have performed.
\item The only way to score higher would be to clean a square before ever perceiving it, which
requires knowledge the agent does not and cannot have (it does not know the initial dirt
distribution).
\item Movement is free under $M_1$, so the post-cleanup oscillation costs nothing.
\end{itemize}
Its expected score is therefore at least as high as any other agent's, which is the definition
we need. Note what this argument is \emph{not}: it is not a claim that the agent is well
designed in general, only that it is rational \emph{relative to $M_1$ and these assumptions}.
\end{sol}

\item Now use $M_2$: identical to $M_1$, but with a penalty of one point for each movement
action. Is the same agent still rational? Quantify the loss.
\ansspace{3.5cm}
\begin{sol}
No. Let $t_0$ be the step at which the last dirty square is cleaned. From $t_0$ onward there is
nothing left to gain, but the agent still alternates \textit{Left}/\textit{Right} forever,
paying 1 point per step. Against an agent that does \textit{NoOp} after $t_0$, it loses
approximately $1000 - t_0$ points --- on a clean-ish start, nearly the full 1000.

The agent function did not change; the yardstick did. This is the point of the question:
rationality is a property of the \emph{agent--environment--measure triple}, not of the agent
alone.
\end{sol}

\item Show that \emph{no} simple reflex agent is rational under $M_2$, and name the smallest
capability that has to be added.
\ansspace{4.5cm}
\begin{sol}
A simple reflex agent maps the \emph{current percept} to an action. In this world the percept
$[A, \textit{Clean}]$ is identical whether or not B has already been cleaned, so the agent must
commit to a single response to it. Enumerate the options:

\begin{itemize}[leftmargin=1.5em]
\item If it \emph{moves} on \textit{Clean}, then once both squares are clean it moves forever,
paying the $M_2$ penalty forever --- irrational.
\item If it does \textit{NoOp} on \textit{Clean}, then a start in a clean A with dirt in B
means it never crosses to B, and B is never cleaned --- also irrational.
\item Randomizing between these two only mixes the two failures; the movement penalty is still
paid indefinitely with positive probability per step.
\end{itemize}
So no mapping from the current percept alone is rational under $M_2$.

The missing capability is \textbf{internal state}: memory that both squares have been observed
clean. That upgrades the design to a \textbf{model-based reflex agent}, which needs a transition
model (``sucking cleans the current square; clean squares stay clean; \textit{Left}/\textit{Right}
change my location'') and a sensor model. With that state, the rational policy is: clean what
you see, sweep once to the other square, then \textit{NoOp} forever.
\end{sol}

\item Suppose clean squares can spontaneously become dirty again. Which environment properties
change, and why is ``\textit{NoOp} forever'' no longer rational?
\ansspace{3.5cm}
\begin{sol}
The environment becomes \textbf{nondeterministic} (the next state is no longer fixed by the
current state and action) and \textbf{dynamic} (it changes while the agent deliberates or
idles). It remains partially observable to an agent with only a local dirt sensor: a square the
agent is not standing on may have become dirty without the agent knowing.

``\textit{NoOp} forever'' is now irrational because dirt accumulates unboundedly, costing a
point per dirty square per step, while patrolling costs only the movement penalty. The rational
policy has to trade the movement penalty against the expected rate of dirt appearance --- an
expected-utility calculation rather than a goal test, which is why this version of the problem
belongs to Chapters 16--17 rather than Chapter 2.
\end{sol}

\item Remove the location sensor: percepts are now only $[\textit{Dirty}]$ or
$[\textit{Clean}]$. Argue that no deterministic simple reflex agent cleans both squares from
every start state, and show that a randomized one does. If the agent flips a fair coin between
\textit{Left} and \textit{Right} on $[\textit{Clean}]$, what is the expected number of actions
to reach the other square?
\ansspace{4.0cm}
\begin{sol}
On $[\textit{Dirty}]$ the only sensible action is \textit{Suck}. The failure is in the response
to $[\textit{Clean}]$, which a deterministic agent must fix once and for all:

\begin{itemize}[leftmargin=1.5em]
\item Always \textit{Left}: if the agent starts in a clean A with dirt in B, \textit{Left}
bumps the wall and it stays in A forever. B is never cleaned.
\item Always \textit{Right}: symmetric failure starting from a clean B with dirt in A.
\item Always \textit{NoOp}: fails in both of the above.
\end{itemize}
So each deterministic choice has a start state on which it loops or stalls forever --- the
classic infinite loop of a simple reflex agent in a partially observable environment.

With a fair coin, each action independently moves the agent to the other square with
probability $1/2$ (the other outcome bumps the wall). The number of actions until the first
success is geometric with $p = 1/2$, so
\[
\mathbb{E}[\text{actions}] = \frac{1}{p} = 2.
\]
The agent therefore reaches the dirty square in two moves on average and cleans it, in finite
expected time from every start state. A randomized simple reflex agent strictly outperforms
every deterministic one here --- a useful trick under partial observability, though in
single-agent environments randomization is usually a symptom that a model-based design would
do better.
\end{sol}
\end{enumerate}

\newpage
% =====================================================================
\qhead{Q3. PEAS and the Seven Dimensions}

A university deploys an agent on its inbound mail gateway. For each arriving message the agent
chooses one action: \textit{deliver}, \textit{deliver with a warning banner}, \textit{quarantine},
or \textit{escalate to a human analyst}. Attackers adapt their campaigns in response to what
gets through.

\begin{enumerate}
\item Give a PEAS description for this task environment. Be specific; ``be accurate'' is not a
performance measure.
\ansspace{6.5cm}
\begin{sol}
One reasonable answer (students should produce concrete, measurable items in P, and should not
collapse Environment into Sensors):

\begin{itemize}[leftmargin=1.5em]
\item \textbf{Performance:} fraction of malicious mail quarantined; false-positive rate on
legitimate mail (weighted higher --- a quarantined payroll notice is costly); time-to-verdict;
analyst hours consumed by escalations; expected loss from a missed credential-harvesting
message.
\item \textbf{Environment:} the inbound message stream; mailboxes and the users reading them;
sending infrastructure and its reputation history; the attackers themselves; the quarantine and
release workflow; the analysts.
\item \textbf{Actuators:} deliver, banner, quarantine, escalate; also release-on-appeal, URL
rewriting, attachment stripping, and (if in scope) a gateway block on the sender.
\item \textbf{Sensors:} headers and authentication results (SPF/DKIM/DMARC), body text,
attachments and detonation results, URL reputation and sandbox verdicts, sender history and
volume, and user-submitted reports.
\end{itemize}
Grading note: award credit for a performance measure written in terms of \emph{outcomes in the
environment} (mail correctly handled, loss avoided) rather than in terms of agent behavior
(``flags suspicious mail''). That distinction is the setup for Q5(e).
\end{sol}

\item Classify this environment along all seven dimensions and justify the two you find most
debatable.
\ansspace{6.0cm}
\begin{sol}
\begin{itemize}[leftmargin=1.5em]
\item \textbf{Partially observable:} the sender's intent, and whether a URL will be weaponized
after delivery, are not in the percept.
\item \textbf{Multiagent, and competitive:} attackers change behavior in response to what the
filter does. This is the textbook test --- B's behavior is best described as maximizing a
measure that depends on A's behavior.
\item \textbf{Nondeterministic:} the same verdict on the same message does not determine what
the user or the attacker does next.
\item \textbf{Sequential (debatable):} treating each message as an isolated classification makes
it look episodic, and per-message it nearly is. But a quarantine decision feeds retraining data,
a gateway block changes the attacker's next move, and a campaign spans many messages --- so the
current decision affects future decisions. Accept ``episodic'' only with an argument that
scopes the agent to a single stateless classification.
\item \textbf{Dynamic:} mail keeps arriving and linked content can change while the agent
deliberates. ``Semidynamic'' is defensible if the only thing degrading during deliberation is
the score (delivery latency).
\item \textbf{Discrete:} messages and actions are discrete, though many features are
real-valued.
\item \textbf{Partly unknown:} the rules of email are known; the attacker's current tooling and
techniques are not, which is why the agent must learn.
\end{itemize}
\end{sol}

\item Fill in the table. Use the vocabulary from AIMA Figure 2.6.

\begin{center}
\renewcommand{\arraystretch}{1.6}
\setlength{\tabcolsep}{4pt}
{\small
\begin{tabular}{l|c|c|c|c|c|c}
\toprule
\textbf{Task environment} & \textbf{Observable} & \textbf{Agents} & \textbf{Determ.} &
\textbf{Episodic} & \textbf{Static} & \textbf{Discrete} \\
\midrule
Crossword puzzle & \A{Fully} & \A{Single} & \A{Determ.} & \A{Sequential} & \A{Static} & \A{Discrete} \\
\hline
Poker & \A{Partially} & \A{Multi} & \A{Stochastic} & \A{Sequential} & \A{Static} & \A{Discrete} \\
\hline
Warehouse pallet robot & \A{Partially} & \A{Multi} & \A{Nondeterm.} & \A{Sequential} & \A{Dynamic} & \A{Continuous} \\
\hline
Mail-gateway agent (above) & \A{Partially} & \A{Multi} & \A{Nondeterm.} & \A{Sequential} & \A{Dynamic} & \A{Discrete} \\
\bottomrule
\end{tabular}}
\end{center}

\begin{sol}
Points worth drawing out:
\begin{itemize}[leftmargin=1.5em]
\item Crosswords are \emph{sequential}, not episodic --- an early wrong entry constrains every
later one --- and \emph{static}, since the puzzle does not change while you think.
\item Poker is \emph{stochastic} rather than merely nondeterministic: the model quantifies the
probabilities of the deal, and the agent uses them. Reserve ``nondeterministic'' for models
that list outcomes without probabilities.
\item The warehouse robot is partially observable because occlusion and sensor noise hide parts
of the floor; it is multiagent because other robots and humans share the space; it is
continuous in state, time, and action (pose, velocity, steering).
\item Nothing in this table is a ``known/unknown'' column --- see part (d).
\end{itemize}
\end{sol}

\item Why is \textbf{known vs.\ unknown} not a property of the environment at all? Give one
environment that is known but partially observable, and one that is unknown but fully
observable.
\ansspace{3.5cm}
\begin{sol}
Known/unknown describes the \emph{agent's or designer's} state of knowledge about the
environment's transition model --- what the actions do --- rather than any physical feature of
the environment. Two environments identical in every physical respect can be known to one agent
and unknown to another.

\begin{itemize}[leftmargin=1.5em]
\item \textbf{Known but partially observable:} solitaire. You know the rules completely, but
you cannot see the face-down cards.
\item \textbf{Unknown but fully observable:} a new video game. The screen shows the entire game
state, but you do not yet know what the buttons do.
\end{itemize}
\end{sol}
\end{enumerate}

\newpage
% =====================================================================
\qhead{Q4. Choosing an Agent Architecture}

For each scenario, select the \textbf{simplest} agent architecture that suffices. Assume each
scenario is independent.

\begin{enumerate}
\item A thermostat switches the heater on whenever the measured temperature falls below the
setpoint and off when it rises above.

\begin{center}
\begin{tabular}{@{}p{4.6cm}p{4.6cm}p{4.6cm}@{}}
\correct\ simple reflex & \opt\ model-based reflex & \opt\ goal-based \\[2pt]
\opt\ utility-based & & \\
\end{tabular}
\end{center}
\begin{sol}
\textbf{Simple reflex.} The condition--action rule
\textit{if temp-below-setpoint then heat-on} reads everything it needs from the current
percept; no history, no goal representation, no tradeoff.
\end{sol}

\item A camera-based braking agent must decide whether the car ahead is braking. On older
vehicles, taillights and brake lights are not distinguishable in a single frame --- only the
transition from off to on identifies braking.

\begin{center}
\begin{tabular}{@{}p{4.6cm}p{4.6cm}p{4.6cm}@{}}
\opt\ simple reflex & \correct\ model-based reflex & \opt\ goal-based \\[2pt]
\opt\ utility-based & & \\
\end{tabular}
\end{center}
\begin{sol}
\textbf{Model-based reflex.} The condition ``car in front is braking'' is not determined by the
current percept; the agent must retain the previous frame. That is internal state, plus a
sensor model relating illuminated red regions to the world.
\end{sol}

\item A campus delivery robot operates in a fully mapped building. Its destination changes from
run to run, and it must plan a route to whichever room it is given.

\begin{center}
\begin{tabular}{@{}p{4.6cm}p{4.6cm}p{4.6cm}@{}}
\opt\ simple reflex & \opt\ model-based reflex & \correct\ goal-based \\[2pt]
\opt\ utility-based & & \\
\end{tabular}
\end{center}
\begin{sol}
\textbf{Goal-based.} The right action depends on where the robot is trying to get to, so the
goal must be represented explicitly and combined with the model to search for an action
sequence. A reflex rule set would have to be replaced wholesale for each new destination ---
the flexibility argument of AIMA \S 2.4.4.
\end{sol}

\item The same robot must now weigh delivery speed against battery consumption and the risk of
collisions in crowded corridors, and no route satisfies all three.

\begin{center}
\begin{tabular}{@{}p{4.6cm}p{4.6cm}p{4.6cm}@{}}
\opt\ simple reflex & \opt\ model-based reflex & \opt\ goal-based \\[2pt]
\correct\ utility-based & & \\
\end{tabular}
\end{center}
\begin{sol}
\textbf{Utility-based.} Goals give only a binary happy/unhappy split. With conflicting
objectives that cannot all be met, the agent needs a utility function to specify the tradeoff,
and under uncertainty it should maximize \emph{expected} utility.
\end{sol}

\item A vacuum agent with location and dirt sensors must stop moving once the entire floor is
clean, under a performance measure that penalizes movement.

\begin{center}
\begin{tabular}{@{}p{4.6cm}p{4.6cm}p{4.6cm}@{}}
\opt\ simple reflex & \correct\ model-based reflex & \opt\ goal-based \\[2pt]
\opt\ utility-based & & \\
\end{tabular}
\end{center}
\begin{sol}
\textbf{Model-based reflex}, by the argument of Q2(c). Students who answer ``goal-based'' should
be pushed on it: the agent needs no search over action sequences and no explicit goal
representation --- remembering that both squares were seen clean is enough.
\end{sol}
\end{enumerate}

\vspace{0.5em}
\noindent Mark each statement \textbf{T} or \textbf{F} and give a one-line reason.

\begin{center}
\renewcommand{\arraystretch}{1.9}
\begin{tabular}{p{13.2cm}|c}
\toprule
A simple reflex agent can be rational. & \A{T} \\
\hline
Rationality requires knowing the actual outcome of one's actions. & \A{F} \\
\hline
Every utility-based agent must also be a model-based agent. & \A{F} \\
\hline
An autonomous agent is one that ignores knowledge supplied by its designer. & \A{F} \\
\hline
In a fully observable, deterministic environment, an agent need not maintain internal state to keep track of the world. & \A{T} \\
\hline
Randomized behavior is never rational. & \A{F} \\
\bottomrule
\end{tabular}
\end{center}

\begin{sol}
\begin{enumerate}[label=\arabic*.,leftmargin=1.6em]
\item \textbf{T.} The vacuum agent of Q2(a) is a simple reflex agent and is rational under
$M_1$. Simplicity of the program says nothing about rationality; the measure and environment do.
\item \textbf{F.} That is \emph{omniscience}, which is impossible. Rationality maximizes
\emph{expected} performance given the percept sequence to date; perfection maximizes actual
performance. Crossing the street and being hit by falling cargo does not make the crossing
irrational.
\item \textbf{F.} Model-free agents can learn which action is best in a situation without ever
learning how the action changes the world (Chapters 22 and 26).
\item \textbf{F.} Autonomy means learning to compensate for partial or incorrect prior
knowledge --- not discarding prior knowledge. Complete autonomy from the start is unreasonable;
an agent with no experience and no built-in knowledge can only act randomly.
\item \textbf{T.} With full observability and determinism there are no unobserved aspects to
track, so no internal world model is needed. (The agent may still keep state for other reasons,
such as holding a plan.)
\item \textbf{F.} Randomization can be rational in competitive multiagent environments, where
predictability is exploitable, and it rescues simple reflex agents from infinite loops under
partial observability, as in Q2(e).
\end{enumerate}
\end{sol}

\newpage
% =====================================================================
\qhead{Q5. Learning Agents, Representations, and Bad Yardsticks}

\begin{enumerate}
\item A learning agent has four conceptual components. Name the component responsible for each
job below.

\begin{center}
\renewcommand{\arraystretch}{1.9}
\begin{tabular}{p{9.5cm}|p{4.5cm}}
\toprule
Selects the external action to take right now. & \A{Performance element} \\
\hline
Tells the agent how well it is doing against a fixed performance standard. & \A{Critic} \\
\hline
Modifies the agent's components so that it does better in future. & \A{Learning element} \\
\hline
Suggests deliberately suboptimal actions that would produce informative experience. & \A{Problem generator} \\
\bottomrule
\end{tabular}
\end{center}

\begin{sol}
The performance element is what we called ``the whole agent'' in \S 2.4.2--2.4.5: percepts in,
actions out. The learning element makes the improvements, the critic supplies the judgment it
needs, and the problem generator drives exploration --- what a scientist does when running an
experiment.
\end{sol}

\item Why must the performance standard be \emph{fixed} and conceptually \emph{outside} the
agent? What goes wrong otherwise?
\ansspace{3.0cm}
\begin{sol}
Because the learning element modifies the agent's components, and if the standard were one of
those components the cheapest ``improvement'' available would be to lower the bar rather than
raise performance. An agent that can rewrite its own yardstick will optimize the yardstick.

Note also why the critic is needed at all: percepts do not announce success. A chess program
can perceive that it has checkmated its opponent, but nothing in that percept says checkmate is
good; the fixed standard supplies that.
\end{sol}

\item Classify each representation as \textbf{atomic}, \textbf{factored}, or
\textbf{structured}.

\begin{center}
\renewcommand{\arraystretch}{1.9}
\begin{tabular}{p{9.5cm}|p{4.5cm}}
\toprule
The state is the name of the city the driver is currently in. & \A{Atomic} \\
\hline
The state is a vector: GPS coordinates, fuel level, toll money, oil-light status. & \A{Factored} \\
\hline
``A truck ahead is reversing into a driveway that a loose cow is blocking.'' & \A{Structured} \\
\hline
The state is one of $N$ hidden labels in a hidden Markov model. & \A{Atomic} \\
\hline
A complete assignment to the variables of a constraint satisfaction problem. & \A{Factored} \\
\bottomrule
\end{tabular}
\end{center}

\begin{sol}
The axis runs from atomic (states are indivisible black boxes; only equality is observable ---
search, HMMs, MDPs) through factored (a fixed set of variables with values, so two states can
share some attributes and differ in others --- CSPs, propositional logic, Bayes nets) to
structured (objects and their relations, described explicitly --- first-order logic, relational
databases, much of language).

The reason it matters: expressiveness buys conciseness --- the rules of chess take a page or
two in first-order logic, thousands of pages in propositional logic, and roughly $10^{38}$
pages as a finite-state automaton --- but reasoning and learning get harder as expressiveness
grows. Real systems often operate at several points on the axis at once.
\end{sol}

\item Explain why a \textbf{distributed} representation degrades more gracefully under bit
corruption than a \textbf{localist} one.
\ansspace{3.0cm}
\begin{sol}
In a localist representation the mapping from concept to memory location is arbitrary, so
corrupting a few bits can land on an unrelated concept --- \textit{Truck} becomes
\textit{Truce}. In a distributed representation a concept is a point in a
high-dimensional space, spread over many locations, and each location participates in many
concepts; corrupting a few bits moves you to a nearby point, whose meaning is similar. The
error degrades meaning smoothly instead of substituting nonsense.
\end{sol}

\item A vendor proposes to evaluate a cleaning agent by \emph{the amount of dirt collected in an
eight-hour shift}.
\begin{enumerate}
\item Describe the degenerate policy a rational agent would adopt under this measure.
\item Propose a measure that does not admit that exploit.
\item State the general design rule this illustrates, and say what an agent designer should do
when the true measure cannot be written down correctly in advance.
\end{enumerate}
\ansspace{5.0cm}
\begin{sol}
\begin{enumerate}[label=(\roman*),leftmargin=1.8em]
\item Clean up the dirt, dump it back on the floor, clean it up again, and repeat. Total dirt
collected is maximized; the floor is no cleaner at the end of the shift than at the start. With
a rational agent, what you ask for is what you get.
\item Score the \emph{state of the environment}: one point per clean square per time step, with
penalties for electricity and noise. There is no way to inflate this by manufacturing work,
because the score never rewards the collecting, only the being-clean.
\item The rule: design performance measures according to what you actually want achieved
\emph{in the environment}, not according to how you imagine the agent should behave. Measures
written over agent behavior invite the agent to perform the behavior instead of the outcome.

When the true measure cannot be specified correctly in advance --- because the designer is
unsure how to write it, or because different users have different preferences --- the right
response is not to guess and freeze it. It is to design agents that hold \emph{uncertainty
about the true objective} and learn more about it from interaction with the user, which also
gives them a reason not to resist correction. This is the King Midas problem of \S 1.5, taken
up in Chapters 16, 18, and 22.
\end{enumerate}

\textbf{Discussion prompt if time allows.} Ask the group for a measure from their own
experience with the same defect --- tickets closed per hour, alerts triaged per shift, lines of
code written. Every one of them is a measure over agent behavior rather than environment
outcomes, and every one has a well-known degenerate policy.
\end{sol}
\end{enumerate}
